\section{Extended Discussion: Reading the Three Findings Together}
\label{sec:discext}
Taken separately, each of this paper's findings has a mundane reading.
A 21.15\% confident-joint disagreement estimate could be an artefact of a weak
probe --- the gate exists to answer that. A +5.9-point cleaning gain could be a seed fluke --- the pre-registered multi-seed replication says worse: random removal matches flagged removal, so we claim no causal reading at all. A 45.7-point head-to-head win could mean the external tool
was never meant for paediatric JPEGs --- the deep dive says exactly
that. The case for the findings is that they \emph{converge}:
\begin{enumerate}
\item The census says the training labels are one-directionally
contaminated (973:23), concentrated on acquisition-degraded films
($p=1.7\times10^{-27}$).
\item Removing candidate flags improves a compact model more than any hyperparameter change we tested (+5.9 points vs +0.4 for Optuna's best) --- but random removal matches it, so this item stands as an unattributed regime observation, not convergent evidence.
\item The failure surface of the strongest external tool is the mirror
image of the same stratum: it misses 384 of 390 pneumonias while never
crying wolf on normals --- a model calibrated away from exactly the
borderline region where this benchmark's labels and images disagree.
\end{enumerate}
Three independent instruments --- a label-process audit, a training
intervention, and an external-tool autopsy --- point at the same
structure. That convergence, not any single number, is the paper's
evidential core.
\subsection{Medical relevance, stated carefully}
This work does not diagnose pneumonia, and we do not frame it clinically
--- that framing would invite validation questions a label-audit study
cannot and should not answer. The relevance is infrastructural: medical
AI benchmarks are the substrate on which diagnostic models are developed,
compared and reproduced, and hidden annotation inconsistencies distort
all three. This collection's labels were generated by automated report
parsing and it anchors a large published literature (150 records
screened here). An audit that publishes per-image, falsifiable,
tiered candidates is a maintainable contribution to that infrastructure:
it lets every downstream user of the benchmark know which films carry
annotation risk, at what confidence, and why.
\subsection{What we would do with more compute}
Ordered by expected information per GPU-hour, updated after the
Round-2 external review: (1) a five-seed consensus-frequency study -
rerun the gated census at five seeds, release per-image flag frequency
(image\_id, times\_flagged/runs, tier) and the Jaccard matrix, turning
the current two-run tiering into the definitive stability artifact;
(2) a consensus-tier removal ablation - remove the 91-image consensus
core and the low-confidence tier separately under the same protocol to
test whether the multi-run tier is the scientifically active object;
(3) a blinded stratified re-read of 50 consensus-core, 50 low-frequency
and 50 control films by a paediatric radiologist (the definitive
validation, and cheap in money rather than compute); (4) recalibrated
head-to-heads against the external tool with its threshold fit on a
validation split; (5) a pretrained-backbones census crosscheck at
higher resolution. None requires new data or new methods; all are
scheduled as future work in the repository.
